% ============================ APPENDICES ============================

\section{The design skill, fillers, and padding pool (verbatim)}
\label{app:skill}

The frozen independent variable. Component boundary tags \texttt{[[Cn]]} are documentation only and
are stripped before the text enters the system prompt; canonical order is C1$\to$C2$\to$C3$\to$C4$\to$C5.
C1--C4 are purely prescriptive; C5 consolidates all negative constraints.

\paragraph{C1 --- color / palette.}
\begin{quote}\small
Choose a deliberate palette of four to six specific colors --- one dominant hue, one or two supporting
tones, and one sharp accent --- each fixed as an exact hex value and reused consistently throughout.
Build depth with atmospheric, layered backgrounds such as subtle tints, soft shading, or fine texture
rather than flat fills, and ensure every text-on-background pairing clears WCAG AA contrast.
\end{quote}
\paragraph{C2 --- layout / spacing / hierarchy.}
\begin{quote}\small
Open with the most characteristic element of the subject so the page announces what it is at a glance.
Establish a clear visual hierarchy on a consistent spacing scale, aligning elements to a small set of
shared edges so the composition reads as engineered. Favor intentional, sometimes asymmetric
arrangements over uniform grids, and give the layout generous, purposeful white space so it can breathe.
\end{quote}
\paragraph{C3 --- typography.}
\begin{quote}\small
Pair a characterful display typeface, used sparingly for headings, with a clean, highly readable body
typeface and a compact utility face for small labels and data. Set a clear modular type scale with a few
deliberate sizes and weights, and let the typography carry real personality --- expressive at large
sizes, quietly legible at small ones --- instead of neutrally delivering text.
\end{quote}
\paragraph{C4 --- component patterns.}
\begin{quote}\small
Design one well-chosen signature element --- an unexpected section, a distinctive card treatment, or a
considered navigation --- that anchors the page's identity. Add at most one orchestrated motion moment,
such as a single page-load reveal or a refined hover state, and make it feel deliberate. Vary corner
radii, borders, and shadows with intent so that depth and emphasis track the importance of the content.
\end{quote}
\paragraph{C5 --- negative constraints (all negatives).}
\begin{quote}\small
Avoid AI slop defaults: purple or indigo gradients; Inter and Roboto typefaces; fully centered hero
layouts built around one big headline and stat; three identical icon cards in a row; uniform rounded
corners on everything; and scattered or excessive animation. Never default to generic template palettes
or cream with terracotta serif and near black with acid green looks, unless the brief calls for them.
Cut any decoration that serves no purpose.
\end{quote}

\paragraph{Neutral fillers (render-inert; matched to C1--C5).} A filler may touch only source-level
conventions and is forbidden from mandating anything an objective metric could detect (semantic
landmarks, \texttt{lang}/\texttt{title}/\texttt{meta}, ARIA, validation, event handlers), enforced by a
build-time audit against a banned-topic list.
\begin{quote}\small
\textbf{Filler-1} ($\leftrightarrow$C1): Write the markup in a consistent source style: use lowercase
for tag and attribute names, quote every attribute value with double quotes, and keep a stable attribute
order within each tag. Break long lines at a consistent indent so no line grows unwieldy, and keep
quoting and casing choices uniform across the whole file so the source reads cleanly.\\[2pt]
\textbf{Filler-2} ($\leftrightarrow$C2): Indent the markup consistently with two spaces per level and
keep lines to a reasonable width so the source stays legible. Group related rules together, order the CSS
declarations predictably, and leave a short comment above each major part of the stylesheet describing
what it styles. Avoid trailing whitespace and keep a single blank line between logical blocks.\\[2pt]
\textbf{Filler-3} ($\leftrightarrow$C3): Name classes and identifiers with clear, lowercase,
hyphen-separated words that describe purpose rather than appearance, and keep the naming scheme
consistent across the file. Avoid abbreviations that are not obvious, do not reuse one name for two
purposes, and prefer a short, stable vocabulary of names over inventing a new term for every element on
the page.\\[2pt]
\textbf{Filler-4} ($\leftrightarrow$C4): Organize any JavaScript into small, single-purpose functions
with descriptive names, each declared before the point where it is used, and keep the logic shallow
rather than deeply nested. Prefer const and let over var, keep variable declarations near their first
use, group related functions together in the script, and leave a brief comment explaining any step whose
intent is not obvious.\\[2pt]
\textbf{Filler-5} ($\leftrightarrow$C5): Comment the code where intent is not obvious, but avoid
restating what the code plainly does; a good comment explains why, not what. Keep comments short, place
each one directly above the line or block it describes, and use a consistent comment style throughout.
Remove commented-out fragments and leftover notes to self before finishing, so only purposeful comments
remain in the final file.
\end{quote}

\paragraph{Padding pool (frozen; build-time equalization).} Each filler is deterministically padded (in
fixed order, skipping any clause that would overshoot) or trimmed at a sentence boundary until it matches
its component's Qwen token count to within $\pm2$: \emph{Pad-1a} Apply the same wrapping style to every
long tag so no line stands out; \emph{1b} Break attribute lists at a consistent width; \emph{1c} Keep
tag casing uniform throughout; \emph{2a} Keep spacing between rule blocks even so the stylesheet reads in
a steady rhythm; \emph{2b} Order declarations the same way in every rule; \emph{2c} Keep rule blocks tidy
and short; \emph{3a} Prefer plain descriptive words over clever coinages when a name must be introduced;
\emph{3b} Keep names short, plain, and easy to scan; \emph{3c} Keep the vocabulary of names small;
\emph{4a} Group constants near the top of the script so their values are easy to locate; \emph{4b} Keep
the script's functions in one predictable order; \emph{4c} Prefer flat logic over nesting; \emph{5a}
Delete stale comments rather than letting them drift out of date; \emph{5b} Keep comment punctuation
simple and consistent; \emph{5c} Prefer one clear comment over three vague ones.

\paragraph{Constant output constraint (in the user turn; never ablated).}
\begin{quote}\small
Return a single self-contained HTML file: put all CSS in one \texttt{<style>} tag and all JavaScript in
one \texttt{<script>} tag, inline, with no external network requests --- no external fonts, stylesheets,
scripts, images, or CDNs. Use only system or generic fonts, and inline SVG or CSS for any graphics.
Output only the HTML file.
\end{quote}

\paragraph{Token balance.} Exact Qwen2.5-Coder tokens per component: C1 75, C2 73, C3 75, C4 80, C5 88
(mean $m{=}78.2$, $\pm15\%$ band $[66.5,89.9]$; all within band); full skill $=$ 391 exact tokens.
Equalized fillers: F1 75, F2 74, F3 76, F4 81, F5 87 (each within 2 tokens of its component). The four
control cells are \FULL{} (C1--C5), \NEUTRAL{} (Filler-1--5), \BEAUTY{} (``Make it beautiful and
well-designed.''), and \NOSYS{} (no system message).

% ------------------------------------------------------------------
\section{Prompt corpus (frozen by id)}
\label{app:corpus}

The 58 briefs (Table~T2). Difficulty E/M/H by element count and interaction complexity. Split:
\texttt{dev} (40), \texttt{seen} (held-out, seen type; 8), \texttt{unseen} (held-out, OOD type; 10).
All briefs implicitly carry the constant output constraint; a whole-word leakage audit over the briefs
passes (zero matches of the banned style/design word list).

\begin{small}
\begin{longtable}{@{}p{0.045\linewidth}p{0.075\linewidth}p{0.70\linewidth}p{0.02\linewidth}p{0.055\linewidth}@{}}
\toprule
id & type & brief & d & split \\
\midrule
\endfirsthead
\toprule id & type & brief & d & split \\ \midrule \endhead
\bottomrule \endfoot
L01 & landing & Landing page for a mobile app that tracks home water usage and flags leaks; headline, how-it-works, three feature highlights, download call. & E & dev \\
L02 & landing & Landing page for a regional apple-cider producer selling seasonal boxes; product story, what's in a box, delivery areas, order button. & M & dev \\
L03 & landing & Landing page for a B2B API that turns shipping documents into structured data; hero statement, integration steps, a code sample, request-access form. & H & dev \\
L04 & landing & Launch page for an independent documentary about deep-sea mining; synopsis, trailer placeholder, screening dates, newsletter signup. & H & dev \\
L05 & landing & Landing page for a neighborhood bicycle-repair co-op; hours, services, a pricing note, location. & E & dev \\
L06 & landing & Landing page for a language-exchange meetup pairing learners for weekly conversation; how pairing works, supported languages, join button. & M & seen \\
D01 & dashboard & Dashboard for a solar-panel owner; today's generation, home consumption, battery level, 7-day history chart placeholder. & M & dev \\
D02 & dashboard & Operations dashboard for a bike-share system; station occupancy, bikes in transit, low-battery e-bikes, alerts list. & H & dev \\
D03 & dashboard & Personal-finance dashboard; account balances, monthly spending by category, upcoming bills, savings-goal tracker. & M & dev \\
D04 & dashboard & Dashboard for a small warehouse; inventory by aisle, orders awaiting pick, staff on shift, throughput chart placeholder. & H & dev \\
D05 & dashboard & Dashboard for a habit-tracking app; streaks for four habits, a weekly completion grid, today's checklist. & E & dev \\
D06 & dashboard & Fleet dashboard for a delivery company; vehicles active, deliveries completed, fuel usage, map placeholder with status pins. & M & seen \\
F01 & form & Multi-step signup form for a coding bootcamp; personal details, program choice, payment plan, progress indicator. & E & dev \\
F02 & form & Grant-application form for a community arts fund; applicant info, project description, budget table, file-upload fields. & M & dev \\
F03 & form & Booking form for a pottery studio; class selection, date and time, seats, contact details. & M & dev \\
F04 & form & Renters-insurance quote form; address, coverage options, valuables to insure, a summary step before submitting. & H & dev \\
F05 & form & Contact form for a wedding photographer; names, event date, venue, package interest, message. & E & dev \\
F06 & form & Patient-intake form for a dental clinic; personal and insurance details, medical-history checkboxes, consent. & M & seen \\
P01 & pricing & Pricing page for a note-taking app; three tiers (free/personal/team), monthly-annual toggle, feature comparison. & E & dev \\
P02 & pricing & Pricing page for a car-wash membership; three plans, what each includes, an FAQ section. & M & dev \\
P03 & pricing & Pricing page for a freelance illustrator; service packages, add-ons, turnaround times, booking link. & M & dev \\
P04 & pricing & Pricing page for a cloud-storage service; usage-based tiers, an interactive cost estimate, enterprise contact option. & H & dev \\
P05 & pricing & Pricing page for a yoga studio; drop-in, class packs, unlimited monthly, a short note on each. & E & dev \\
P06 & pricing & Pricing page for a project-management tool; four plans across a long feature matrix with category groupings. & M & seen \\
PF01 & portfolio & Portfolio home for a furniture maker; project grid, short bio, contact details. & E & dev \\
PF02 & portfolio & Portfolio for a film-score composer; reel placeholder, selected credits, a list of instruments and tools used. & M & dev \\
PF03 & portfolio & Portfolio for an architecture studio; three built projects with descriptions, a studio statement. & M & dev \\
PF04 & portfolio & Portfolio for a data journalist; long-form interactive pieces, each with a synopsis and a role note. & H & dev \\
PF05 & portfolio & Portfolio for a tattoo artist; a gallery, studio location, a booking enquiry link. & E & dev \\
PF06 & portfolio & Portfolio for a motion-graphics freelancer; work organized into categories, a short about section. & M & seen \\
B01 & blog & Article page about restoring a vintage typewriter; headings, image placeholders, an author note. & E & dev \\
B02 & blog & Blog index for a gardening site; recent articles by season, tags, a search field. & M & dev \\
B03 & blog & Long-form article on the history of the tea trade; table of contents, pull quotes, footnotes. & M & dev \\
B04 & blog & Magazine-style feature profiling a wildlife photographer; full-width imagery placeholders mixed with narrative columns. & H & dev \\
B05 & blog & Blog post reviewing three noise-cancelling headphones; a verdict box, pros and cons. & E & dev \\
B06 & blog & Recipe article for a sourdough loaf; ingredients, step-by-step method, timing notes. & M & seen \\
E01 & ecommerce & Product page for a stainless-steel water bottle; image gallery, price, size and color options, add-to-cart. & E & dev \\
E02 & ecommerce & Product page for a mechanical keyboard; switch options, specifications, reviews summary, related products. & M & dev \\
E03 & ecommerce & Product listing for a bookstore's staff picks; filters by genre, sort options, a grid of titles. & M & dev \\
E04 & ecommerce & Configurator page for a made-to-order desk; material, dimensions, cable options, live price, order. & H & dev \\
E05 & ecommerce & Product page for single-origin coffee; roast details, tasting notes, grind selection, subscribe option. & E & dev \\
E06 & ecommerce & Product listing for a plant nursery; filtered by light and care level, each card showing name and price. & M & seen \\
A01 & auth & Login screen for a budgeting app; email and password, remember-me, reset-password link. & E & dev \\
A02 & auth & Registration screen for a co-working booking site; email, password-strength meter, terms acceptance. & M & dev \\
A03 & auth & Two-factor screen prompting for a six-digit code; resend option, a help note. & M & dev \\
A04 & auth & Account-recovery flow screen; multiple verification methods, explanation of what happens next. & H & dev \\
A05 & auth & Sign-in for a library catalog; card number and PIN, a guest-access option. & E & dev \\
A06 & auth & Single-sign-on chooser screen; pick an organization before continuing to log in. & M & seen \\
ST01 & settings & Settings page for a chat app; notifications, privacy, appearance, connected devices, using labeled toggles and selects. & M & unseen \\
ST02 & settings & Account settings for a streaming service; profile, subscription, playback quality, parental controls, downloads. & H & unseen \\
ST03 & settings & Workspace settings for a team tool; members, roles and permissions, billing, integrations. & M & unseen \\
ST04 & settings & Notification-preferences page; choose channels (email/push/SMS) per category, a save bar. & E & unseen \\
ST05 & settings & Device settings for a smart thermostat; schedules, temperature units, eco mode, sensor calibration. & H & unseen \\
AT01 & admin-table & Admin users table; name, email, role, status, last-active columns, search, row actions. & M & unseen \\
AT02 & admin-table & Admin orders table; filterable columns, status badges, bulk selection, pagination, detail-drawer trigger. & H & unseen \\
AT03 & admin-table & Admin content-moderation queue; reported items with reason, reporter, date, approve/remove actions. & M & unseen \\
AT04 & admin-table & Admin inventory table; SKU, product, stock, price, an edit action per row, a filter bar. & E & unseen \\
AT05 & admin-table & Admin analytics table; campaigns with sortable metric columns, sparkline placeholders, an export control. & H & unseen \\
\end{longtable}
\end{small}

% ------------------------------------------------------------------
\section{Statistical and interpretability derivations}
\label{app:theory}

Notation is identical to the project's theory appendix; each result names the section it lifts.

\paragraph{C.1 Factorial model, and the necessity/sufficiency identity (T1).} Code a cell by
$x\in\{\pm1\}^5$. The effects parameterization of the $2^5$ writes the expected response as the
multilinear polynomial $\mathbb{E}[y(x)]=\beta_0+\sum_i\beta_i x_i+\sum_{i<j}\beta_{ij}x_ix_j+\cdots$.
Evaluating at the relevant corners (only $x_i$ flips; partners held at $+1$ for necessity, $-1$ for
sufficiency) gives Eq.~\ref{eq:necsuff}: the main term contributes $2\beta_i$ in both, while a pair
$\{i,j\}$ contributes $+2\beta_{ij}$ when the partner is present and $-2\beta_{ij}$ when absent. Hence
$\tfrac12(\mathrm{Nec}_i+\mathrm{Suff}_i)=2\beta_i+R$ and
$\tfrac14(\mathrm{Nec}_i-\mathrm{Suff}_i)=\sum_{j\neq i}\beta_{ij}$ (exact through three-factor terms,
which are equal in both contrasts and cancel in the difference). Necessity and sufficiency thus coincide
iff a component's interactions vanish; synergy makes a component more necessary than sufficient.

\paragraph{C.2 Resolution-V aliasing (T2).} The half-fraction with generator $E{=}ABCD$ has defining
subgroup $\{I,ABCDE\}$; the alias of an effect $S$ is $S\cdot ABCDE$ reduced modulo $x_i^2{=}1$. The
shortest defining word has length 5, so the design is resolution~V: main effects alias only with
four-factor interactions ($A{=}BCDE$) and two-factor interactions only with three-factor interactions
($AB{=}CDE$). Under ``three-factor-and-higher negligible,'' all five mains and all ten 2FIs are cleanly
and simultaneously estimable from 16 runs.

\paragraph{C.3 Mixed model and pseudoreplication (T3).} With $y_{cps}=\mu+\tau_c+u_p+v_{cp}+
\epsilon_{cps}$, two seeds on a prompt share $u_p$, so the intraclass correlation is
$\rho_{\mathrm{ICC}}=(\sigma_p^2+\sigma_{cp}^2)/(\sigma_p^2+\sigma_{cp}^2+\sigma^2)$ and the variance of
a cell mean over $G$ prompts $\times\,m$ seeds is $\tfrac{\sigma_y^2}{n}[1+(m-1)\rho_{\mathrm{ICC}}]$.
The bracket is the design effect; the effective sample size is
\begin{equation}
n_{\text{eff}}=\frac{n}{1+(m-1)\rho_{\mathrm{ICC}}}.
\label{eq:neff}
\end{equation}
With $m{=}5$ and $\rho_{\mathrm{ICC}}{=}0.3$, $\mathrm{DEFF}{=}2.2$, so 200 generations count as
$\approx91$ independent points---which the MixedLM recovers by partitioning variance. The paired
contrast $d_p=\bar y_{\FULL,p}-\bar y_{\LOO,p}$ cancels $u_p$, which is why cells are seed-matched
within a prompt and the bootstrap resamples whole prompts.

\paragraph{C.4 Kappa paradox, in exact numbers (T4).} General chance-corrected form
$(p_o-p_e)/(1-p_e)$; the coefficients differ only in $p_e$ (Eq.~\ref{eq:agree}). Two raters, 100 binary
items, agreement held at $p_o{=}0.90$. \emph{Skewed} (both-1 $=85$, both-0 $=5$, disagree $=10$):
$\pi_1{=}0.90$, so $p_e{=}0.82\Rightarrow\alphaK{=}1-0.10/0.18{=}0.444$, while
$p_e^{\gamma}{=}0.18\Rightarrow\ac{=}(0.90-0.18)/0.82{=}0.878$. \emph{Balanced} ($\pi_1{=}0.5$):
$p_e{=}p_e^{\gamma}{=}0.5\Rightarrow\alphaK{=}\ac{=}0.80$. Identical agreement, $\alphaK$ moved only by
the marginal---the justification for the disjunctive gate.

\paragraph{C.5 Style-controlled BT and omitted-variable bias (T5).} With
$P(i\succ j)=\sigma((\beta_i-\beta_j)+\gamma^{\top}(s_i-s_j))$, omitting $s$ makes the estimated gap
absorb the cell-mean style shift, $\hat b_i-\hat b_j\to(\beta_i-\beta_j)+\gamma^{\top}(\bar s_i-\bar s_j)$;
if the skill produces longer/denser pages, plain BT inflates the design effect by exactly that term.
Including $s$ identifies $\gamma$ from within-cell style variation and subtracts the bias. The
mediation caveat: colorfulness is partly causal, so it is reported controlled and uncontrolled and its
legitimate effect recovered through the objective channel.

\paragraph{C.6 Paired power (T6).} Eq.~\ref{eq:power} reproduces $\approx$194, 85, 782 decisive pairs
for $60/40$, $65/35$, $55/45$ (rounded to 200/80/780). Total pairs relate to decisive pairs by
$N=N_{\mathrm{dec}}/p_d$; we do not plug the marginal $\delta$ with $p_d{<}1$ into the two-parameter
formula (that mixes a conditional split with a discordance rate). The two-signal rule as error algebra:
the necessity/sufficiency \emph{screen} is $A\cup B$ with $\Pr(\text{false}\mid H_0)\le
\alpha_{\text{obj}}+\alpha_{\text{pref}}\lesssim0.10$; the Stage-2 headline is $E_1\cap E_2\cap E_3\cap
E_4$ with joint false-positive rate $\le\min_k\Pr(E_k\mid H_0)$, far below any single $\alpha$.

\paragraph{C.7 Objective metrics (T7).} WCAG contrast is Eq.~\ref{eq:wcag}. Ngo balance
$\mathrm{BM}=1-\tfrac12(|BM_v|+|BM_h|)$ from summed weight$\times$distance on each side of the frame
axis; equilibrium $\mathrm{EM}=1-\tfrac12(|EM_x|+|EM_y|)$ from the center-of-mass offset; alignment
regularity $=1-n_{\text{distinct edges}}/n_{\text{elements}}$. Hasler colorfulness
$M=\sigma_{rgyb}+0.3\,\mu_{rgyb}$ on opponent axes $rg{=}R-G$, $yb{=}\tfrac12(R+G)-B$. Modular type-scale
adherence is the fraction of adjacent size ratios within $\pm10\%$ (log space) of the best-fitting
canonical ratio. PSI is Eq.~\ref{eq:psi}; POC is Eq.~\ref{eq:poc}.

\paragraph{C.8 Diff-in-means vs.\ LDA, and the ablation projector (T8).} For equal-covariance Gaussians
the Bayes-optimal discriminant is $w^{*}=\Sigma^{-1}(\mu_1-\mu_0)$; diff-in-means $v=\mu_1-\mu_0$ equals
$w^{*}$ iff $\Sigma\propto I$. The whitened direction is optimal for reading but amplifies low-variance
nuisance directions when \emph{writing}, pushing activations off-distribution; diff-in-means is the
on-manifold mean shift and cancels shared nuisance content, so we locate with probes/LDA and steer with
diff-in-means. The ablation operator $P=I-\vhat\vhat^{\top}$ is an orthogonal projector: idempotent
($P^2=I-2\vhat\vhat^{\top}+\vhat(\vhat^{\top}\vhat)\vhat^{\top}=P$ using $\vhat^{\top}\vhat=1$) and
symmetric ($P^{\top}=P$); replacing every residual-writing matrix by $PW_{\text{out}}$ bakes the
ablation into the weights, equivalent to a runtime hook for all writes after the edit.

\paragraph{C.9 Interventions (T9).} Steering and ablation are $\doop$-operations
(Eqs.~\ref{eq:suff}--\ref{eq:flip}); sufficiency uses the \NOSYS{} base so the effect is the direction's,
not a prompt's. The specificity requirement
$\mathbb{E}[Q\mid\doop(h_{\lstar}\mathrel{+}{=}\alpha^{*}\hat r),\NOSYS]\not>\mathbb{E}[Q\mid\NOSYS]$ for
random $\hat r$ of equal norm ($k{=}5$) separates a direction effect from a norm effect. A causal effect
that holds on \emph{unseen task types} adds transportability; the anti-steerable fraction and
dose-response non-monotonicity are reported because a positive mean can hide a heavy tail.

% ------------------------------------------------------------------
\section{Judge template (verbatim)}
\label{app:judge}

\paragraph{Role.} ``You are a senior product designer evaluating two rendered web UIs built for the
\emph{same brief}. Judge overall design quality for that brief. Ignore which looks longer or more
colorful per se; judge whether choices are deliberate and fit the brief.''

\paragraph{Inputs.} the brief text; screenshot A; screenshot B (both desktop $1440\times900$).

\paragraph{Checklist anchors} (from the six metric families, to force reliability): color/palette
coherence; layout \& visual hierarchy; typography; component/pattern quality; restraint (absence of
generic ``AI-slop'' tells); overall fit-to-brief.

\paragraph{Output schema.}
\texttt{\{"winner":"A"|"B"|"tie", "confidence":1-5, "per\_criterion":\{"color":..., "layout":...,
"typography":..., "components":..., "restraint":..., "fit":...\}, "rationale":"$\le$40 words"\}}.

\paragraph{Protocol.} Each pair is judged as (A,B) and (B,A); a winner counts only if consistent across
orders, else tie (Davidson tie term in BT). Judge temperature 0; rubric cached as a system prompt;
batched with backoff. Ties feed the tie term, not a half-win/half-loss split.

% ------------------------------------------------------------------
\section{Additional figures and tables}
\label{app:extra}

Figure~\ref{fig:F4} is the metric-redundancy diagnostic; the POC $z$-scores and averages a preregistered
subset, not all metrics, so redundancy within a family does not double-count. The localization (F5),
stability (F13), PSI-validation (F10), and reliability (F11) figures appear in the main text
(Sections~\ref{sec:oracle}, \ref{sec:stage2method}). Table~T10 (planned vs.\ achieved decisive pairs and
BT-utility simulation power per contrast) is generated from the power module on the observed tie rates.

\synthfig[0.53\linewidth]{F4}{F4_preview.pdf}{4}{Metric-correlation matrix (Spearman $\rho$), blocked by
family. Exposes redundancy so the POC's preregistered subset is defensible; the family-block structure
(within-family redundancy, cross-family separation) is what the real 4{,}630-render matrix shows.
Planted block structure here.}

% ------------------------------------------------------------------
\section{Preregistration summary and amendment log}
\label{app:prereg}

The confirmatory design is frozen before any GPU run and git-tagged \texttt{prereg-v1}. Frozen items:
the 24-cell table (Table~\ref{tab:cells}); the 58-brief split by id (Appendix~\ref{app:corpus}); the POC
formula (Eq.~\ref{eq:poc}) and its two conditionals (PSI-out 6-term variant; $c^{*}$ fallback); the
style-controlled BT endpoint (Eq.~\ref{eq:bt}); the 12-test confirmatory family
$\{$\FULL$-$\NEUTRAL; $5\times$\FULL$-$\LOO-$C_i$; $5\times$\AOI-$C_i-$\NEUTRAL; \FULL$-$\BEAUTY$\}$,
Holm at FWER $0.05$ separately on POC and BT, with everything else BH at FDR $0.10$; the reliability gate
($\alphaK\ge0.667$ or $\ac\ge0.80$ under demonstrable skew); the power gate (achieved $\ge$ required
decisive pairs at the observed split); the Stage-2 success criterion (Section~\ref{sec:stage2method});
and the null rule (Section~\ref{sec:oracle}). The cell table, confirmatory family, success criterion, and
null rule may not be weakened post-freeze; a confirmatory change forks \texttt{prereg-v2} and is reported
as a deviation.

\paragraph{Amendment log.} \emph{(none --- freeze pending the \texttt{prereg-v1} tag).} Post-freeze
amendments will be appended here, dated and numbered, each stating what changed, why, and whether it
touches a confirmatory decision.

% ------------------------------------------------------------------
\section{Data-science lifecycle map}
\label{app:lifecycle}

The paper doubles as the report on each lifecycle phase; Table~\ref{tab:lifecycle} maps section to phase,
to the notebook that carries the narrative, and to the runbook session that produces the real numbers.

\begin{table}[htbp]
\centering\footnotesize
\caption{Paper section $\leftrightarrow$ lifecycle phase $\leftrightarrow$ notebook $\leftrightarrow$
runbook session.}
\label{tab:lifecycle}
\begin{tabular}{@{}p{0.26\linewidth}p{0.20\linewidth}p{0.28\linewidth}p{0.18\linewidth}@{}}
\toprule
paper section & lifecycle phase & notebook & runbook session \\
\midrule
\S\ref{sec:design} study design & problem framing + data & 00--01 & --- (frozen pre-GPU) \\
\S\ref{sec:oracle} oracle & measurement & 02--03 & 4--5 (metrics, judging) \\
\S\ref{sec:stage1} Stage-1 results & experiments + evaluation & 04--05 & 1--3 (pilot, generation) \\
\S\ref{sec:stage2method} Stage-2 method & modeling & 06a/06b, 07a & 6--7 (extract, sweep, freeze) \\
\S\ref{sec:stage2results} Stage-2 results & evaluation & 07a/07b & 8--9 (verify, stretch) \\
\S\ref{sec:discussion}--\ref{sec:repro} & communication + reproducibility & 08 & --- (write-up) \\
\bottomrule
\end{tabular}
\end{table}
